\chapter{Conclusion}
\label{chapter:conclusion}

The Municipality of Utrecht asked how it can programme the maintenance,
strengthening and replacement of more than 60\,km of quay walls when the
evidence is incomplete and the theoretical end of life is a poor guide. This
report proposes a method that ranks walls by risk rather than by age, and tests
it on the walmuren of the Wervengebied, the one area with enough drawings, dive
records and renewal data to build and check a failure model.

The method has four steps. It first checks the asset register, whose errors
decide which walls are looked at: in the test case, 11 of the 29 walmuren that
the register shows as past their design life had already been renewed on steel
piles. It then computes a failure probability where a structural model exists,
here for the 17 walls (611\,m) that still stand on their original timber piles
with a werf, with the evidence that every one of them stands
(Chapter~\ref{chapter:probabilistic-modelling}); walls without a model are
ranked with indicators such as satellite movement and condition scores. It
combines the probability, used only relatively, with a consequence score into
a red, orange or green class (Section~\ref{sec:approach-selection}). Finally it
compares the measures for each class in a fault tree and a decision tree and
extrapolates the pace of renewal (Chapters~\ref{chapter:risk-analysis} and
\ref{chapter:prediction-recommendations}).

Five conclusions follow for the test case. First, the evidence that matters
most is the wall type, which changes $\beta$ by about 2, and the condition of
the piles; both are cheap to establish compared with a renewal of about
\euro1.3 million per wall. Second, the failure model is conservative for the
older wall type and its absolute values are no safety verdict; as a ranking it
makes nine of the 17 walls (398\,m) red, four of them because data are
missing; four walls (274\,m) stay red in every variant tested. Third, at the pace of the last eight
years the 17 walls are renewed around 2031, and about one of them is expected to
reach a limit state before then; if failures are spread over time, the pace of
renewal matters more than its order, but if the weakness is present today,
about three walls may reach it whatever the pace. Fourth, because the walls are renewed within a few years, measures that
act at once are the effective ones: restricting the use of the werf near the
edge is almost free in expected value and the best choice under a pessimistic
view, and renewing the weakest walls early with extra capacity is preferred by
a risk-averse owner; whether restricting use beats the current programme
depends on the absolute failure probability, which the model does not yet give. Monitoring pays for walls that will wait much longer, which
is most of the portfolio. Fifth, several improvements need decisions rather
than money: a register that records renewals and foundations, a consequence
class for each wall type, and acceptance of proven strength.

The results rest on two schematic cross-sections from the Oudegracht, literature
values for the timber and the soil, and assumed probabilities for most human and
organisational events, so they hold as a ranking and an order of magnitude. The
next steps are to establish the wall type and the pile condition of the red
walls, to obtain the municipal unit costs and incident records, and to apply
the indicator score to the rest of the 60\,km, so that the classes of the test
case can be calibrated and the same method can guide the whole programme.
